%% notes-tex/025-additive-baselines/sections/1-question.tex

\section{The question, and the state of the evidence\secstatus{todo}}
\label{sec:question}

Tran et al.\ introduced MULTI-evolve in \emph{Science} in 2026
(\texttt{10.1126/science.aea1820}), a protein engineering framework whose fully
connected networks were said to learn the epistatic landscape from single and
double mutants and to identify synergistic combinations at higher order. Visani,
Verma and DeWitt replied in a preprint, \emph{Additive baselines furnish no evidence
for epistasis learning by MULTI-evolve} (bioRxiv \texttt{2026.04.23.719915}), which
does one thing: it fits the null that the original paper never fit. An
\term{additive model} over mutations predicts a phenotype as a sum of per-mutation
terms and assigns exactly zero to every interaction, so it is the natural null for
any claim that a model has learned interaction, as an intercept-only model is the
null for a claim that a model has learned anything. Their networks' predictions
correlated with a ridge-regularized additive model's at $r > 0.999$, and on held-out
data the network did not significantly outperform the additive model or a
zero-hidden-layer version of itself.

The analogous question for the 010 trigenic transformer is whether its held-out
performance on the trigenic interaction score survives the same test, and it has a
second part the protein setting does not have: the published split lets every model
score the identity of the screen a record came from. Both are answered here on the
025 build, which holds the 010 records with bit-identical labels, under two splits:
the published random split (arm R) and a split that never shares a Kuzmin query
double between parts (arm Q).

\notesec{One difference in the setup, and it cuts against the null}
In the protein setting the quantity predicted is the multimutant phenotype itself,
which per-mutation effects explain directly. Here the quantity predicted is already
a residual: the trigenic interaction score subtracts the expectation built from
single and double mutant fitness before the model ever sees it. An
additive-in-single-gene-effects model is therefore a harder null here than it is
there, and its performance is more surprising, not less. That argument has one
loophole, the screen structure of Section~\ref{sec:data}, and quantifying it turns
out to matter more than anything else below.

\begin{keybox}
\textbf{Measured.} Five nulls and the nonlinear baseline on both arms, on test
(Table~\ref{tab:heldout}, Figure~\ref{fig:ladders}a,b). Arm R reproduces the 010
numbers. On arm Q the additive ridge falls from 0.400 to 0.185, the screen-only
model to zero, and the nonlinear baseline from 0.432 to 0.141, below the ridge.
Four transformer runs on arm Q, every one peaking on validation by epoch 7 and
declining afterward (Figure~\ref{fig:disjointruns}a,b). One transformer test score
on arm Q: GH 1640's validation-selected checkpoint at 0.139, below the ridge by
0.046 with a paired interval that excludes zero, level with the nonlinear baseline
(Table~\ref{tab:decision}).

\textbf{Not measured.} Replicates of the 010 configuration on arm Q, which is
what turns one checkpoint's number into a statement about the model, and the
learnable-table control for the constant-learning-rate runs.
\end{keybox}

\notesec{Vocabulary}
\term{Arm R} and \term{arm Q} are the two training arms of experiment 025, one per
split; every baseline is fit under both. The \term{ladder} is the ordered set of
baselines B0 through B5 and the transformer, read left to right by interaction
capacity. A \term{query pair} is the double mutant a Kuzmin trigenic screen was
built around; every record in the build is one query pair plus one \term{array
gene}, and 420 query pairs recur across the 376{,}732 records. A \term{screen} is
the set of records sharing a query pair. A \term{validation maximum} is the largest
validation Pearson over a run's logged epochs; it is an upward-biased order
statistic, and it is never set beside a test number without that label.
